\documentclass{article}
\usepackage[utf8]{inputenc}
\usepackage{amsmath,amssymb}
\usepackage{hyperref}
\usepackage{graphicx}
\title{Autonomous Discovery in Domain: Calculus}
\author{Cognitive Engine v7}
\date{\today}
\begin{document}
\maketitle
\begin{abstract}
This paper presents an autonomously generated mathematical synthesis discovered by the v7 Cognitive Engine. Notably, the engine successfully escaped local minima by discovering and injecting \textbf{1} novel meta-axioms into the Prolog unified physics ruleset.\end{abstract}
\section{Initial Conditions}
The evolutionary process initialized with the following paradox or seed topology:
\begin{verbatim}
(DERIV (SIN X T) X)
\end{verbatim}
\section{Discovered Meta-Axioms}
To resolve topological deadlocks, the Lisp orchestrator synthesized the following Prolog rules and successfully verified them against the domain invariants:
\begin{itemize}
  \item \texttt{unified_eval(qnm_resonance(M, L), P) :- unified_eval(add(M, L), P)}
\end{itemize}
\section{Optimal Hybrid Topology}
Following the Open-Ended Novelty Search and final Crossover Synthesis, the engine converged on the following optimal topology (Penalty = 0.0):
\begin{verbatim}
(+ (INT X T) (* (/ X T) X))
\end{verbatim}
\section{Computational Metrics}
\begin{itemize}
  \item \textbf{Evaluated ASTs:} 118
  \item \textbf{Elapsed Time:} 0.5792 seconds
  \item \textbf{Throughput:} 203.73 ASTs/sec
\end{itemize}
\end{document}
